\section{Adjoint jets and the saturated semigroup}
\label{adj:section}

We retain the orbifold and jet data of Proposition~\ref{jet:setup}.
In particular, $B$ is a smooth effective proper complex orbifold of
dimension $N$, its coarse space is projective, $L$ is an ample orbifold
line bundle, and $-K_B=rL$.  At the chosen point the stabilizer is
$\mu_m$, the fiber character of $L$ is faithful, and equivariant
transverse parameters $z_1,\ldots,z_N$ have characters
$b_1,\ldots,b_N$.  We use the original, unscaled degree coordinate:
\[
 \Lambda=\left\{(j,\gamma)\in\ZZ\times\ZZ^N:
       j\equiv\sum_i b_i\gamma_i\pmod m\right\},
 \qquad c_0=(r,\mathbf1).
\]
The semigroup $\Gamma$ records the total-degree-then-lexicographic
Taylor initials of sections of $jL$, and $K$ is its closed real cone.
It has full dimension, contains the positive constant ray, and has
bounded slices of bounded original degree.  Its projection onto the
exponent coordinates is the nonnegative orthant.  The canonical
character identity gives $c_0\in\Lambda$.

The adjoint lifting statement below uses these data but does \emph{not}
assume that $c_0$ itself lies in the interior of $K$.  We use
Theorem~\ref{jet:interior} only after proving that statement.

\subsection{Vanishing on the smooth orbifold models}

We first record the precise form of vanishing needed on a weighted
blowup and its resolution.  Divisors and line bundles on a stack are
written additively.  Nefness and bigness of their rational classes are
understood on the projective coarse space.

\begin{lemma}[Orbifold vanishing]\label{adj:vanishing}
Let $\mathcal X$ be a smooth proper Deligne--Mumford stack over $\CC$,
with trivial generic stabilizer and projective coarse space.  Let
$\Theta$ be an effective rational divisor with simple normal crossing
support and coefficients less than one.  If $D$ is a Cartier divisor
and $D-(K_{\mathcal X}+\Theta)$ is nef and big, then
\[
 H^q(\mathcal X,\cO_{\mathcal X}(D))=0\qquad(q>0).
\]
\end{lemma}

\begin{proof}
Let $q:\mathcal X\to X$ be the coarse morphism.  For a prime divisor
$B\subset X$, let $e_B$ be the stabilizer order at the generic point
of its inverse image.  A rational divisor $A$ on $\mathcal X$ has a
coarse rational divisor $A_c$, characterized in codimension one by
$q^*A_c=A$.  Locally the coarse map is given in the normal direction
by $t=z^{e_B}$, so
\[
 K_{\mathcal X}=q^*(K_X+R),
 \qquad R=\sum_B(1-1/e_B)B.
\]
Only divisors with $e_B>1$ contribute to $R$.

The coefficient of $D_c$ at $B$ is an integral multiple of $1/e_B$.
Consequently
\begin{equation}\label{adj:coarse-boundary}
 \Delta=R+\Theta_c-\{D_c\}
\end{equation}
is effective: the fractional coefficient being subtracted is at most
$(e_B-1)/e_B$.  The pair $(X,\Delta)$ is klt.  Indeed, on a smooth
finite quotient uniformizer $U\to X$, the log pullback of $\Delta$
is the pullback of $\Theta$ minus the effective pullback of
$\{D_c\}$.  This boundary is bounded above by the simple normal
crossing boundary $\Theta|_U$ with coefficients less than one.
All its log discrepancies are therefore positive.  The finite-map
discrepancy formula, with ramification included in the displayed
canonical identity, gives the same positivity downstairs.

Normal extension from codimension one identifies
\[
 q_*\cO_{\mathcal X}(D)=\cO_X(\lfloor D_c\rfloor).
\]
For completeness, on a finite quotient chart the pushforward is the
invariant part of an invertible module on a normal smooth scheme.
Sections extend across sets of codimension at least two there and
then remain invariant, so this pushforward is the rank-one reflexive
sheaf prescribed by its codimension-one orders.  Those orders are
exactly $\lfloor D_c\rfloor$.  Taking invariants is exact in
characteristic zero, which also gives
\[
 H^q(\mathcal X,\cO_{\mathcal X}(D))
   =H^q(X,\cO_X(\lfloor D_c\rfloor)).
\]

The coarse space has quotient singularities and is locally
$\QQ$-factorial.  Thus $\lfloor D_c\rfloor$ is an integral
$\QQ$-Cartier Weil divisor.  By Equation~\eqref{adj:coarse-boundary},
\[
 \lfloor D_c\rfloor-(K_X+\Delta)
   =D_c-(K_X+R+\Theta_c),
\]
which is the nef and big class descended from
$D-(K_{\mathcal X}+\Theta)$.  Kawamata--Viehweg vanishing for a klt
pair and an integral $\QQ$-Cartier Weil divisor now applies
\cite{FujinoKV}.  It gives the claimed vanishing.  This argument
requires neither flatness of the coarse morphism nor ampleness in
place of nefness and bigness.
\end{proof}

\subsection{The adjoint lifting implication}

\begin{lemma}[Adjoint lattice points]\label{adj:implication}
For the orbifold jet data above,
\begin{equation}\label{adj:adjoint-implication}
 (j,\alpha)\in\Lambda,\qquad
 (j,\alpha)+c_0\in\Int K
 \quad\Longrightarrow\quad (j,\alpha)\in\Gamma.
\end{equation}
\end{lemma}

\begin{proof}
The coordinate inequalities of $K$ imply
$\alpha_i+1>0$.  Since $\alpha_i$ is integral, $\alpha_i\ge0$.
The bounded-slice property and the positivity of the degree coordinate
on $K\setminus\{0\}$ give $j+r>0$.  Choose a rational number
\[
 0<j_0<j+r
 \quad\hbox{such that}\quad
 p=(j_0,\alpha+\mathbf1)\in\Int K.
\]
We construct a section of $jL$ with the required Taylor initial.

\paragraph{Replacing finitely many initials by scalar weights.}
There are finitely many sections whose initial indices span a
polyhedral cone containing $p$ in its interior.  To see this without
assuming finite generation of $\Gamma$, choose a small simplex around
$p$ inside $K$ and approximate its vertices by points in the cone
generated by $\Gamma$.  Sufficiently close approximations still
surround $p$, and each uses only finitely many semigroup elements.
The union of those finite lists has the desired property.

For this list of sections, choose positive primitive integral weights
$u_1,\ldots,u_N$ such that scalar $u$-order has exactly their specified
Taylor initials, each as a single monomial.  Here is the finite-order
justification.  Let $d$ exceed the total degrees of those initials
and the total degree of $\alpha$.  Weights sufficiently close to a
common positive value preserve all strict total-degree comparisons
through degree $d$ and place every monomial of total degree greater
than $d$ above the relevant initials.  Among monomials of equal
degree at most $d$, successive sufficiently small rational
perturbations of the weights realize the chosen lexicographic order.
There are only finitely many such comparisons.  We can include all
monomials preceding $\alpha$ among them, and then clear denominators
and divide by the common divisor.  Thus, in addition,
\begin{equation}\label{adj:scalar-order-choice}
 \gamma\hbox{ precedes }\alpha
       \quad\Longrightarrow\quad u\cdot\gamma<u\cdot\alpha.
\end{equation}
This handles formal Taylor series as well as polynomial expansions,
since positivity of the weights excludes all sufficiently high total
degrees at once.

Set
\[
 p_0=u\cdot\alpha,
 \qquad U=u\cdot(\alpha+\mathbf1)=p_0+\sum_i u_i.
\]
The affine plane given by original degree $j_0$ and scalar cost $U$
meets the finite cone just chosen in a relative neighborhood of $p$.
Choose finitely many rational points in that neighborhood surrounding
$p$ in the plane.  Their expressions as nonnegative combinations of
the integral cone generators can be taken rational.  Clear all
denominators simultaneously.  Multiplication of the corresponding
sections then gives a positive integer $a$ and sections
$s_1,\ldots,s_t$ of $aj_0L$ with single $u$-initials
$z^{\beta^{(1)}},\ldots,z^{\beta^{(t)}}$, after rescaling, such that
\begin{equation}\label{adj:surrounding}
 u\cdot\beta^{(k)}=aU,
 \qquad
 \alpha+\mathbf1\in
 \operatorname{relint}\operatorname{conv}
       \{\beta^{(1)}/a,\ldots,\beta^{(t)}/a\}
       \subset\{x:u\cdot x=U\}.
\end{equation}
The relative interior here contains a neighborhood in the entire
displayed affine plane.  Products preserve single initials, so no
claim about generation of the full section ring has entered this
construction.

\paragraph{The weighted blowup stack.}
Use the algebraic equivariant parameters fixed in
Section~\ref{jet:section}.  For $k\geq0$, let $I_k$ be the ideal
generated locally by monomials of $u$-weight at least $k$.
Each $I_k$ contains an ordinary power of the maximal ideal, so its
quotient is determined on a finite infinitesimal neighborhood of the
chosen residual gerbe.  Equivariance descends this quotient to $B$,
and extending its kernel as the unit ideal off the center gives a
global coherent ideal.  On a chart with a larger finite group, use
the translates at all points of the finite orbit.  One fixed parameter
system defines every $I_k$, so these ideals form a multiplicative
filtration, including after passage to the completed Taylor ring.

The weighted Rees algebra
\[
 \mathcal A=\bigoplus_{k\geq0}I_k\mathsf q^k
\]
is locally generated over the coordinate ring by
$z_i\mathsf q^\ell$, $1\leq\ell\leq u_i$: distribute $k$ among
the factors of a monomial of weight at least $k$, assigning at most
$u_i$ to each occurrence of $z_i$.  Define
$\pi:\mathcal Z\to B$ as the stack quotient by the grading
$\Gm$ of $\Spec_B\mathcal A\setminus V(\mathcal A_+)$.
This global construction commutes with the etale parameter charts
and is the identity away from the center, where $I_k=\cO_B$.

We identify its local smooth model to retain the exceptional inertia.
Put $x_i=z_i\mathsf q^{u_i}$.  These elements cover the irrelevant
open by their nonvanishing loci, because
\[
 (z_i\mathsf q^\ell)^{u_i}=x_i^\ell z_i^{u_i-\ell}.
\]
On $x_i\ne0$, the ratio
$t=(z_i\mathsf q^{u_i-1})/x_i$ is regular, including when $u_i=1$,
and $z_j=t^{u_j}x_j$.  Thus before taking the original cyclic quotient
our model is
\[
 \left[
  \bigl(\A^{N+1}_{x_1,\ldots,x_N,t}
       \setminus\{x_1=\cdots=x_N=0\}\bigr)/\Gm
 \right],
 \qquad
 \rho\cdot(x_i,t)=(\rho^{u_i}x_i,\rho^{-1}t).
\]
The original $\mu_m$ acts by its characters on $x_i$ and trivially
on $t$.  The atlas is smooth and all stabilizers are finite.
The exceptional Cartier divisor $E$ is $t=0$; it is
$[\PP(u_1,\ldots,u_N)/\mu_m]$, with any generic stabilizer retained.
The coarse modification is the relative projective weighted Rees
blowup.  A sufficiently divisible $\cO_{\mathcal Z}(-qE)$ descends
to its relatively ample tautological line bundle.  Since the coarse
space of $B$ is projective, so is that of $\mathcal Z$.

The Jacobian calculation in these coordinates gives
\begin{equation}\label{adj:weighted-canonical}
 K_{\mathcal Z}+E=\pi^*K_B+\Bigl(\sum_i u_i\Bigr)E.
\end{equation}
Indeed $\bigwedge_i dz_i$ has order $\sum_i u_i-1$ along $t=0$ on
each uniformizer.  Keeping the stack structures makes this the
Cartier divisor identity asserted in
Equation~\eqref{adj:weighted-canonical}.

Define the line bundle divisors
\[
 P_0=\pi^*(jL)-p_0E,
 \qquad H_0=\pi^*(aj_0L)-aUE.
\]
The canonical shift cancels the exceptional order because
$U=p_0+\sum_i u_i$.  Together with $-K_B=rL$, this gives
\begin{equation}\label{adj:positive-difference}
 P_0-(K_{\mathcal Z}+E)-\frac1aH_0
       =(j+r-j_0)\pi^*L.
\end{equation}
Its right side is nef and big, with strictly positive coefficient;
this is the positivity needed to lift from $E$.
After their common exceptional order is removed, the sections
$s_1,\ldots,s_t$ define sections of $H_0$.  Let $\mathfrak b$ be
their base ideal.  It does not vanish identically on $E$, and its
restriction there is generated locally by the indicated monomial
initials.

\paragraph{Multiplier-ideal membership on the exceptional divisor.}
The congruence $(j,\alpha)\in\Lambda$ says that $z^\alpha$ has
exactly the character required to be a global section of $P_0|_E$.
Its scalar degree is $p_0$, and $\cO_E(-E)=\cO_E(1)$ in the
weighted projective convention.  We claim that
\begin{equation}\label{adj:multiplier-membership}
 z^\alpha\in H^0\left(E,P_0|_E\otimes
                 \mathcal J_E((\mathfrak b|_E)^{1/a})\right).
\end{equation}

On the chart where $x_k\ne0$, pass to the smooth finite uniformizer
that sets $x_k=1$.  More precisely the chart is
$[\A^{N-1}/H_k]$, where
\[
 H_k=\{(\rho,\zeta)\in\Gm\times\mu_m:
                    \rho^{u_k}\zeta^{b_k}=1\}.
\]
This is a finite extension of $\mu_m$ by $\mu_{u_k}$, acting on the
remaining coordinates by
$x_i\mapsto\rho^{u_i}\zeta^{b_i}x_i$.  It need not be a direct
product or act effectively.  On the underlying affine uniformizer,
$P_0|_E$ and $H_0|_E$ are trivial line bundles with the respective
$H_k$-characters $\rho^{p_0}\zeta^j$ and
$\rho^{aU}\zeta^{aj_0}$.  Thus the generators of the restricted
ideal become the ordinary monomials
$x_{\rm rem}^{\beta^{(i)}_{\rm rem}}$, and the proposed section
becomes $x_{\rm rem}^{\alpha_{\rm rem}}$.  For example, on $H_k$,
\[
 \prod_{i\ne k}(\rho^{u_i}\zeta^{b_i})^{\alpha_i}
       =\rho^{p_0}\zeta^j,
\]
by the defining relation of $H_k$ and the congruence for $\alpha$.
There are no fractional exponents: the root needed to set $x_k=1$
has already been taken in the smooth atlas.  Any generic stabilizer
of $E$ is retained in $H_k$ and imposes exactly these character
conditions, without changing the smooth-atlas multiplier formula.
Projection from the plane $u\cdot x=U$ by
dropping coordinate $k$ is an affine isomorphism.  Hence
Equation~\eqref{adj:surrounding} places
$\alpha_{\rm rem}+\mathbf1$ in the ordinary interior of the convex
hull of the projected scaled generators, and therefore in the
interior of their scaled Newton polyhedron.

The monomial multiplier-ideal formula says that $x^\nu$ belongs to
$\mathcal J(I^c)$ precisely when $\nu+\mathbf1$ is in the interior
of $c$ times the Newton polyhedron of $I$ \cite{Howald01}.  Its
hypotheses hold on this smooth affine chart.  Applying it with
$c=1/a$ proves the claimed membership on every uniformizer.  The
ideal and the section are equivariant, so these memberships descend
and prove Equation~\eqref{adj:multiplier-membership}.  No ordinary
smoothness of the coarse exceptional divisor is being assumed.

\paragraph{Adjoint vanishing and restriction.}
We apply the positivity in Equation~\eqref{adj:positive-difference}
on a projective log resolution
$\sigma:\mathcal Y\to\mathcal Z$ of $E$ and $\mathfrak b$,
isomorphic at the generic point of $E$.  We use smooth stacks
throughout: functorial log resolution in smooth etale charts
\cite[Theorems~35--36]{KollarResolution07}
descends.  Arrange that the total boundary, including
the strict transform $\widetilde E$, has simple normal crossings,
and that the restricted morphism
$\tau:\widetilde E\to E$ resolves $\mathfrak b|_E$.  Write
\[
 \mathfrak b\cO_{\mathcal Y}=\cO_{\mathcal Y}(-G_0).
\]
The effective divisor $G_0$ does not contain $\widetilde E$, and
$\sigma^*H_0-G_0$ is globally generated by the pulled-back sections.

Consider the integral line bundle divisor
\[
 A_{\mathcal Y}=K_{\mathcal Y}
       +\sigma^*(P_0-K_{\mathcal Z}-E)-\lfloor G_0/a\rfloor.
\]
Its difference from $K_{\mathcal Y}+\{G_0/a\}$ is
\begin{equation}\label{adj:resolution-positivity}
 (j+r-j_0)\sigma^*\pi^*L
          +\frac1a(\sigma^*H_0-G_0).
\end{equation}
The first term is nef and big and the second is nef.  The fractional
boundary has simple normal crossings and coefficients less than
one.  Lemma~\ref{adj:vanishing} gives
$H^1(\mathcal Y,\cO_{\mathcal Y}(A_{\mathcal Y}))=0$.
The divisor restriction sequence therefore gives a surjection
\begin{equation}\label{adj:restriction-surjection}
 H^0(\mathcal Y,A_{\mathcal Y}+\widetilde E)
    \longrightarrow
 H^0(\widetilde E,(A_{\mathcal Y}+\widetilde E)|_{\widetilde E}).
\end{equation}

Adjunction identifies the bundle on the right with
\begin{equation}\label{adj:restricted-adjoint}
 \tau^*(P_0|_E)+K_{\widetilde E}-\tau^*K_E
               -\left\lfloor (G_0|_{\widetilde E})/a\right\rfloor.
\end{equation}
Here round-down commutes with restriction.  Indeed, simple normal
crossings of $\widetilde E+\Supp G_0$ prevents distinct components
of $G_0$ from sharing a divisor on $\widetilde E$.  If an individual
intersection is disconnected, its coefficient is the same on each
component, so the equality still holds.  Moreover
$G_0|_{\widetilde E}$ is the divisor of the pulled-back restricted
ideal.  Thus Equation~\eqref{adj:restricted-adjoint} is precisely
the bundle in the log-resolution definition of the multiplier ideal
appearing in Equation~\eqref{adj:multiplier-membership}.
That membership supplies a section of
Equation~\eqref{adj:restricted-adjoint} whose rational section is
the pullback of $z^\alpha$.  Surjectivity in
Equation~\eqref{adj:restriction-surjection} lifts it to
$\mathcal Y$.

\paragraph{Regular pushdown and recovery of the initial.}
The difference between the lifting bundle and $\sigma^*P_0$ is
\begin{equation}\label{adj:pushdown-correction}
 K_{\mathcal Y}-\sigma^*K_{\mathcal Z}
       -\sigma^*E+\widetilde E-\lfloor G_0/a\rfloor.
\end{equation}
The part preceding the last term is effective and exceptional over
$\mathcal Z$.  This can be checked during a resolution by smooth
blowups: a center of codimension $c\ge2$ contributes $c-1$, or
$c-2$ when contained in the strict transform of $E$, in addition
to pullbacks of already effective exceptional terms.  The same
calculation holds in the etale smooth charts.  The final term in
Equation~\eqref{adj:pushdown-correction} is the negative of an
effective divisor.  It follows that the lifted rational section has
no pole along any nonexceptional prime divisor.  It pushes down to
a regular section of $P_0$ on the normal stack $\mathcal Z$.
Its restriction to $E$ agrees generically with $z^\alpha$, and hence
agrees everywhere, since both are regular sections on the integral
smooth stack $E$.

Finally $p_0\ge0$, so $P_0$ is a subsheaf of $\pi^*(jL)$.
Normal extension across the center, or the local weighted blowup
charts, gives $\pi_*\cO_{\mathcal Z}=\cO_B$.  Pushing down thus
produces a section of $jL$.  The exceptional order and the specified
restriction say exactly that its $u$-initial is $z^\alpha$, with no
term of smaller $u$-cost.  Equation~\eqref{adj:scalar-order-choice}
then excludes every Taylor monomial preceding $\alpha$ in the
original order.  This proves Equation~\eqref{adj:adjoint-implication}.
\end{proof}

\subsection{Integral support forms and saturation}

We now use the interior conclusion from the preceding section.
The next argument explains why the adjoint implication forces a
polyhedral cone, rather than merely imposing a condition on facets
that might happen to be rational.

\begin{theorem}[The saturated jet semigroup]\label{adj:saturation}
Suppose $c_0\in\Int K$.  Then the following statements hold:
\begin{enumerate}
 \item $K$ is a rational polyhedral cone.
 \item Every facet has a primitive nonnegative support form
       $\eta\in\Lambda^*$ satisfying $\eta(c_0)=1$.
 \item The semigroup and its interior lattice points satisfy
 \begin{equation}\label{adj:saturation-identities}
   \Gamma=K\cap\Lambda,
   \qquad \Int K\cap\Lambda=c_0+\Gamma.
 \end{equation}
\end{enumerate}
In particular, $\Gamma$ is a finitely generated saturated affine
semigroup with group $\Lambda$.  These conclusions apply to every
sufficiently far grading and chosen point covered by
Theorem~\ref{jet:interior}, including its exact rational case.
\end{theorem}

\begin{proof}
Fix such jet data, with $c_0\in\Int K\cap\Lambda$.  Put
$\varphi(x)=(r/N)f(x)$, so the epigraph of $\varphi$ uses the
original degree coordinate $j$, rather than $D_0=Nj/r$.
On the open positive orthant, $\varphi$ is a finite convex function
and is locally Lipschitz.

\paragraph{A support at a differentiability point has height one.}
Let $x$ be a differentiability point of $\varphi$ in the positive
orthant, and let $y=(\varphi(x),x)\in\partial K$.  Homogeneity
gives $\varphi(x)=\nabla\varphi(x)\cdot x$.  The unique supporting
normal ray at $y$ is therefore generated by
\[
 \eta(j',x')=j'-\nabla\varphi(x)\cdot x',
\]
and $\eta(c_0)>0$.  Suppose either that this ray is irrational
relative to $\Lambda$, or that its primitive integral generator
has value greater than one at $c_0$.  In both cases there exists
$w\in\Lambda$ with
\begin{equation}\label{adj:forbidden-strip}
 -\eta(c_0)<\eta(w)<0.
\end{equation}
For an integral primitive form of height at least two, take an
element of value $-1$.  For an irrational ray, the subgroup
$\eta(\Lambda)\subset\RR$ is dense: any nondense additive subgroup
of $\RR$ is discrete, and discreteness here would make a positive
multiple of $\eta$ primitive integral.  This proves the assertion
in the irrational case as well.

There are integers $k_\nu\to\infty$ and lattice points
$l_\nu\in\Lambda$ such that
\begin{equation}\label{adj:recurrence}
 e_\nu=l_\nu-k_\nu y\longrightarrow0.
\end{equation}
This is recurrence of the multiples of $y$ in the compact torus
$\RR^{N+1}/\Lambda$.  It also follows directly by the pigeonhole
principle applied to successively finer finite partitions of a
lattice fundamental domain.  If the resulting return times stay
bounded, one of them is an exact period and its unbounded multiples
give Equation~\eqref{adj:recurrence} instead.

Now $l_\nu+w$ is outside $K$ for large $\nu$, since its
$\eta$-value tends to the negative number $\eta(w)$.  In contrast,
\begin{equation}\label{adj:inward-recurrence}
 l_\nu+w+c_0
   =k_\nu\left(y+\frac{w+c_0+e_\nu}{k_\nu}\right)
      \in\Int K
\end{equation}
for large $\nu$.  To justify this strict inclusion, set
$v=w+c_0$, so $\eta(v)>0$ by
Equation~\eqref{adj:forbidden-strip}.  Differentiability gives,
uniformly for $v'$ in a small neighborhood of $v$,
\[
 \varphi(x+t v'_x)
  =\varphi(x)+t\nabla\varphi(x)\cdot v'_x+o(t).
\]
Thus the $j$-coordinate of $y+t v'$ exceeds
$\varphi(x+t v'_x)$ for every sufficiently small positive $t$;
the exponent coordinates remain positive as well.  Apply this with
$t=1/k_\nu$ and $v'=v+e_\nu$ to obtain
Equation~\eqref{adj:inward-recurrence}.
Lemma~\ref{adj:implication}, applied to the lattice point
$l_\nu+w$, would then put this point in $\Gamma\subset K$, a
contradiction.  Every supporting ray at a differentiability point
therefore has a primitive integral form of height one at $c_0$.

\paragraph{Only finitely many support forms occur.}
Choose $\epsilon>0$ such that the Euclidean ball of radius
$\epsilon$ about $c_0$ is contained in $K$.  If a nonnegative
support form has $\eta(c_0)=1$, testing it on this ball gives
$\|\eta\|\le1/\epsilon$.  Since $\Lambda^*$ is a lattice, only
finitely many primitive integral forms satisfy this bound.  Denote
the forms obtained at differentiability points by
$\eta_1,\ldots,\eta_t$.

These forms, together with the coordinate inequalities, define
$K$.  Indeed, convex functions are differentiable almost everywhere
on their open domain.  At any given positive $x_0$, take
differentiability points $x_\nu\to x_0$.  Local Lipschitz continuity
bounds their gradients, and therefore
\[
 \varphi(x_\nu)+\nabla\varphi(x_\nu)\cdot(x_0-x_\nu)
       \longrightarrow\varphi(x_0).
\]
Each expression is the value at $x_0$ of a supporting linear
function.  Hence the supports at these differentiability points
recover the epigraph on the positive orthant.  Since only the finite
list above occurs, $K$ and
\[
 K'=\{(j,x):x_i\ge0\text{ for all }i,
                     \ \eta_a(j,x)\ge0\text{ for all }a\}
\]
agree there.  We already have $K\subset K'$.  Conversely, for
$q\in K'$, the point $q+t c_0$ lies in $K'$ and has positive
exponent coordinates whenever $t>0$.  It is thus in $K$, and
closedness gives $q\in K$ on letting $t\downarrow0$.  This proves
$K=K'$ and rational polyhedrality.

The coordinate form $(j,x)\mapsto x_i$ is primitive on $\Lambda$:
one can prescribe $x=e_i$ and choose $j$ in the requisite congruence
class.  Its value on $c_0$ is one.  All other defining forms also
have height one.  Every facet of the full-dimensional polyhedral
cone is cut out by a form from an irredundant sublist, so its
primitive support has the asserted height.

\paragraph{Saturation and the interior translation.}
If $q\in K\cap\Lambda$, then $q+c_0\in\Int K$: translate a
small ball about $c_0$ by $q$ and use addition in the convex cone.
Lemma~\ref{adj:implication} gives $q\in\Gamma$, proving the first
identity in Equation~\eqref{adj:saturation-identities}.
If $q\in\Int K\cap\Lambda$, each primitive facet form takes a
positive integral value on $q$, hence a value at least one.  Since
its value on $c_0$ is one, all facet inequalities hold for $q-c_0$.
Thus $q-c_0\in K\cap\Lambda=\Gamma$.  The reverse inclusion
$c_0+\Gamma\subset\Int K\cap\Lambda$ follows from the same
ball-translation argument.  This proves the second identity.

Finally, triangulate the rational polyhedral cone into finitely many
rational simplicial cones and choose lattice generators for their
rays.  A lattice point in any such cone is a nonnegative integral
combination of its ray generators plus a lattice point in their
bounded half-open parallelepiped.  There are only finitely many
such remainder points.  Taking their union over the triangulation
proves finite generation of $K\cap\Lambda$.  If a positive integral
multiple of $q\in\Lambda$ lies in this semigroup, then $q\in K$
by homogeneity, so $q$ already lies in it; this is saturation.
The group is $\Lambda$ by Proposition~\ref{jet:setup}.
\end{proof}
